\section{Obserwacja osoby podczas przejścia}
Celem PATH1 jest uzyskanie kilku kierunków obserwacji osoby, nie skan biometryczny ani rekonstrukcja jej całej powierzchni. Drogę przyjęto w osi $x=20$ m od D01 na $y=0$ do $y=9$ m w R102. Otwarty pas ma 1,8 m szerokości i nie zawiera bramek obrotowych ani przesłon. Bezpośrednie wyjście prowadzi przez D01 lub E03; obserwacja nie wymusza postoju ani dodatkowego odcinka ewakuacji.

Wymaganym ujęciem jest górna część osoby: barki, górny tułów i głowa w bezwzględnym zakresie wysokości 1,2--2,0 m. Nie przyjęto widoczności stóp ani całej sylwetki w tych czterech kadrach. Obwiednia osoby w rzucie ma 0,6 $\times$ 0,6 m; środek może odchylić się o $\pm0{,}6$ m od osi drogi, pozostając w wolnym pasie 1,8 m. Dla osób, których głowa leży niżej niż 1,2 m, nie deklaruje się równoważnego ujęcia. Brak takiego kadru nie blokuje dostępu uprawnionej osoby ani wyjścia.

\PathViewTable

C01 patrzy z północy na przód osoby idącej na północ, C02 z południa na tył, C03 z zachodu na lewy profil, C04 ze wschodu na prawy profil. Wszystkie cztery odcinki mają wspólną część $y=4{,}2$--5,5 m. Podane minima px/m obejmują także krańce odcinków, odchylenie środka, głębokość obwiedni oraz oba poziomy wysokości. Współrzędne obrazu są znormalizowane; każdy wynik mieści się w $\pm0{,}96$ i w granicach near/far.

Porównano przebieg prosty 9 m i załamany przez $(20,4)$, $(22,4)$, $(22,9)$, długości 11 m. Zakręt zmienia ustawienie ciała, ale bez dodatkowych kamer nie zapewnia tyłu i obu profili. Przesłony na zakręcie zasłaniałyby część kadru, a ruch boczny utrudniałby mijanie i przejazd wózka. Ponieważ cztery stałe kamery uzyskują wymagane kierunki na prostej, wybrano prosty przebieg.

Kontrola geometrii używa prostokątnego kadru 3D, HFOV/VFOV, wysokości, pochylenia i granic promieniowych. Obliczono narożniki obwiedni na obu krańcach odcinka oraz w krokach do 0,1 m wewnątrz R102. Nie ma ścian ani szaf między tymi kamerami a osią przejścia. Skrzydło D02 i wyposażenie pozostają poza pasem $x=19{,}1$--20,9 m. W modelu zapisano także regały i szafę IT w innych pomieszczeniach, aby nie rysować kadrów przez ściany lub meble.

Inna osoba, niesiony bagaż lub zwrot głowy mogą zasłonić obraz; sam kształt drogi nie jest dowodem widoczności każdej strony. Przyjęty odstęp operacyjny 3 s odpowiada teoretycznie 1200 osobom/h, lecz do doboru przyjęto tylko 240 osób/h wobec zapotrzebowania 120 osób/h. To założenia przepustowości do próby, nie wyniki obserwacji ruchu. Próba obejmuje grupy, bagaż, osoby niższe, wózki i obu kierunki ruchu; niedostępny kadr nie blokuje wyjścia.
